% GENERATED FILE. Edit canonical lessons, then run npm run generate:books.
% canonical-source-hash: fnv1a64:c19d4438c06c313b
% canonical-lessons: HI-C74-open, HI-C74-closed, HI-C74-entrance, HI-C74-exit, HI-C74-late

\chapter{Words on the Door}
\label{ch:hi-words-on-the-door}
\hlchaptermodality{81}

\begin{chapteropening}
\textbf{By the end of this chapter:} \emph{I can read whether a place is open or closed, find the entrance and the exit, and understand that a train is late before I take my leave.}

Complete HI-C74-late and say all five words in order, then say that the train is late.
\end{chapteropening}

\section[khulā]{\hi{खुला} (khulā) — open}

\label{lesson:HI-C74-open}



\begin{quote}
Two words back: what does \emph{veṭar} mean, and what is \textbf{\hi{पीना}}? The last five words of this run are the ones painted on doors.
\end{quote}

\subsection*{You'll want to know: \hi{खुला}}
\textbf{\hi{खुला}} (\emph{khulā}) — \textquotedblleft{}open\textquotedblright{}.

It is built from \textbf{\hi{खुलना}}, \textquotedblleft{}to open\textquotedblright{}, and it names the state that opening leaves behind — the way English \emph{opened} stands to \emph{open}. Hindi makes these already-done forms constantly, and this is one you will meet before you are taught the pattern.

\hi{खुला} is the word on a shutter, a shop board, a gate and a ticket window. It is also used of a thing that is loose or unpacked: \textbf{\hi{खुला} \hi{पैसा}} is loose change, coins rather than a note, which is the phrase that will save you at a counter.

The \hi{ु} under the \hi{ख} is the short \emph{u} matra, hanging below.

\begin{grammarlens}[title={What you've built}]
The first of five words you read rather than say.
\end{grammarlens}

\subsection*{Your turn}
\begin{itemize}
\raggedright
  \item \emph{Recall:} say \emph{veṭar}, then read \textbf{\hi{पीना}} and say what it means
  \item \emph{Read it:} \hi{खुला}, as though it were painted on a shutter
  \item \emph{Say it:} \emph{khulā}
  \item \emph{Say it:} \emph{khulā paisā}
\end{itemize}

\begin{tcolorbox}[breakable,colback=teal!4,colframe=teal!35!black,title={Before you move on}]
What does \hi{खुला} mean? (\textquotedblleft{}open\textquotedblright{}.) Say it once more.
\end{tcolorbox}
\section[band]{\hi{बंद} (band) — closed}

\label{lesson:HI-C74-closed}



\begin{quote}
Say \emph{khulā}, then write \textbf{\hi{खुला}}: \hi{ख} with \hi{ु} beneath it, then \hi{ल}, then \hi{ा}.
\end{quote}

\subsection*{You'll want to know: \hi{बंद}}
\textbf{\hi{बंद}} (\emph{band}) — \textquotedblleft{}closed\textquotedblright{}.

Persian \ar{بند} (\emph{band}), \textquotedblleft{}tied, shut\textquotedblright{}, and its family is enormous. The ancient root meant \textquotedblleft{}to bind\textquotedblright{}, and English \emph{bind}, \emph{band} and \emph{bond} are its western cousins. A closed door is a \textbf{tied} door.

Two things to keep straight. First, \hi{बंद} is not the beginning of \textbf{\hi{बंदर}}, the monkey you met among the animals — the eye will want to join them and they are unrelated. Second, \hi{बंद} stands opposite \hi{खुला} on exactly the same signboards, so learn them as one pair and you have read half a street.

\begin{grammarlens}[title={What you've built}]
The pair every shutter in the country is painted with.
\end{grammarlens}

\subsection*{Your turn}
\begin{itemize}
\raggedright
  \item \emph{Recall:} read \textbf{\hi{पीना}}, then say it without looking
  \item \emph{Say it:} \emph{khulā}, then \emph{band}
  \item \emph{Read it:} \hi{बंद}, and check that it is not \hi{बंदर}
  \item \emph{Write it:} \hi{खुला} once more, then \hi{बंद} beside it
\end{itemize}

\begin{tcolorbox}[breakable,colback=teal!4,colframe=teal!35!black,title={Before you move on}]
What does \hi{बंद} mean? (\textquotedblleft{}closed\textquotedblright{}.) And its opposite? (\emph{khulā}.)
\end{tcolorbox}
\section[pravesh]{\hi{प्रवेश} (pravesh) — entrance}

\label{lesson:HI-C74-entrance}



\begin{quote}
Say the pair a shutter is painted with, in both orders.
\end{quote}

\subsection*{You'll want to know: \hi{प्रवेश}}
\textbf{\hi{प्रवेश}} (\emph{pravesh}) — \textquotedblleft{}entrance\textquotedblright{}.

Two Sanskrit pieces. \textbf{\hi{प्र}-} is \textquotedblleft{}forward\textquotedblright{}, and \textbf{\hi{विश्}} (\emph{viś-}) is \textquotedblleft{}to enter\textquotedblright{}. Going forward into a place: \hi{प्रवेश}.

That \hi{प्र}- is worth keeping. It is the same ancient piece as Latin \emph{pro-} and English \emph{fore-}, and once you can see it you will spot it at the front of a great many formal Hindi words.

This is a \textbf{written} register, not a spoken one. Nobody says \hi{प्रवेश} to a friend at a doorway; it is painted over the door, printed on a form, and read out in a station announcement. That is the same shelf \textbf{\hi{कृपया}} sits on — the please of signs rather than the please of counters.

The opening \hi{प्र} stacks \hi{प} on \hi{र} with the vowel-killer, as \hi{क्र} did in \hi{क्रिकेट}.

\begin{grammarlens}[title={What you've built}]
Three: \hi{खुला}, \hi{बंद}, \hi{प्रवेश}. Two states of a door and one of its two labels.
\end{grammarlens}

\subsection*{Your turn}
\begin{itemize}
\raggedright
  \item \emph{Read it:} \hi{प्रवेश}, and find the stacked \hi{प्र}
  \item \emph{Say it:} \emph{pravesh}
  \item \emph{Say it:} \emph{pravesh band hai}
\end{itemize}

\begin{tcolorbox}[breakable,colback=teal!4,colframe=teal!35!black,title={Before you move on}]
What does \hi{प्रवेश} mean? (\textquotedblleft{}entrance\textquotedblright{}.) Where do you meet it, on a sign or in speech? (On a sign.)
\end{tcolorbox}
\section[nikās]{\hi{निकास} (nikās) — exit}

\label{lesson:HI-C74-exit}



\begin{quote}
Write \textbf{\hi{प्रवेश}} from memory, starting with the stacked \hi{प्र}.
\end{quote}

\subsection*{You'll want to know: \hi{निकास}}
\textbf{\hi{निकास}} (\emph{nikās}) — \textquotedblleft{}exit\textquotedblright{}.

You have most of this word already. \textbf{\hi{निकलना}}, \textquotedblleft{}to set out, to go out\textquotedblright{}, has been yours since the chapter on leaving; \hi{निकास} is what that doing word turns into when you need a noun. Going out, named.

Watch the vowel change as the noun is made: \hi{निक}\textbf{\hi{ल}}\hi{ना} has a short \emph{a}, \hi{निक}\textbf{\hi{ा}}\hi{स} takes the long \hi{ा}. Hindi lengthens a vowel to build a name out of an action often enough that you should expect it rather than be surprised.

\hi{निकास} is painted opposite \hi{प्रवेश} over every second doorway in a station, and the two are read together or not at all.

\begin{grammarlens}[title={What you've built}]
Four, and the second pair of this chapter: \hi{प्रवेश} and \hi{निकास}.
\end{grammarlens}

\subsection*{Your turn}
\begin{itemize}
\raggedright
  \item \emph{Say it:} \emph{nikalnā}, then \emph{nikās}
  \item \emph{Say it:} \emph{pravesh}, then \emph{nikās}
  \item \emph{Write it:} \hi{निकास}, lengthening the vowel that \hi{निकलना} keeps short
\end{itemize}

\begin{tcolorbox}[breakable,colback=teal!4,colframe=teal!35!black,title={Before you move on}]
What does \hi{निकास} mean? (\textquotedblleft{}exit\textquotedblright{}.) And \hi{प्रवेश}? (\textquotedblleft{}entrance\textquotedblright{}.)
\end{tcolorbox}
\section[der]{\hi{देर} (der) — lateness, a delay}

\label{lesson:HI-C74-late}



\begin{quote}
Name the four words of this chapter so far, in their two pairs.
\end{quote}

\subsection*{You'll want to know: \hi{देर}}
\textbf{\hi{देर}} (\emph{der}) — \textquotedblleft{}lateness, a delay\textquotedblright{}.

Persian \ar{دیر} (\emph{der}), \textquotedblleft{}late, a long while\textquotedblright{}. It is a naming word rather than a describing one: Hindi does not say \textquotedblleft{}the train is late\textquotedblright{}, it says the train \textbf{is with lateness} — \textbf{\hi{गाड़ी} \hi{देर} \hi{से} \hi{है}}.

Set it against \textbf{\hi{जल्दी}}, which you already own for speed and for earliness. \hi{जल्दी} and \hi{देर} are the two ends of the same measure, and a station announcement will use one or the other of them about your train.

That is where this chapter lands. \textbf{\hi{गाड़ी} \hi{देर} \hi{से} \hi{है}} is a sentence built entirely from words you now hold, and it is the single most useful thing a traveller can understand on a platform.

\begin{grammarlens}[title={What you've built}]
Five: \hi{खुला}, \hi{बंद}, \hi{प्रवेश}, \hi{निकास}, \hi{देर}. Enough to read a door, find a way through it, and understand why the train has not come.
\end{grammarlens}

\subsection*{Your turn}
\begin{itemize}
\raggedright
  \item \emph{Say it:} \emph{der}
  \item \emph{Say it:} \emph{gāṛī der se hai}
  \item \emph{Contrast:} \emph{jaldī} and \emph{der}
  \item \emph{Write it:} \hi{निकास}, then \hi{देर}, and say all five in order
\end{itemize}

\begin{tcolorbox}[breakable,colback=teal!4,colframe=teal!35!black,title={Before you move on}]
How do you say the train is late? (\emph{gāṛī der se hai}.) Say it once more.
\end{tcolorbox}
